% 苏布拉马尼扬·钱德拉塞卡（综述）
% license CCBYSA3
% type Wiki

本文根据 CC-BY-SA 协议转载翻译自维基百科\href{https://en.wikipedia.org/wiki/Subrahmanyan_Chandrasekhar}{相关文章}

\begin{figure}[ht]
\centering
\includegraphics[width=6cm]{./figures/4962eacafc5afbdc.png}
\caption{} \label{fig_Sbl_1}
\end{figure}
苏布拉马尼扬·钱德拉塞卡（Subrahmanyan Chandrasekhar，/ˌtʃəndrəˈʃeɪkər/，音译：昌德拉谢卡；[4] 泰米尔语：சுப்பிரமணியன் சந்திரசேகர்，罗马化：Cuppiramaṇiyaṉ Cantiracēkar；1910年10月19日－1995年8月21日）[5] 是一位印裔美国理论物理学家，在恒星结构、恒星演化以及黑洞等领域作出了卓越贡献。他在职业生涯中也曾将相当一部分时间投入于流体力学研究，特别是稳定性与湍流理论，并在这些领域取得了重要成果。钱德拉塞卡与威廉·A·福勒因在“与恒星结构与演化有关的重要物理过程的理论研究”方面的贡献，共同获得1983年诺贝尔物理学奖。他对恒星演化的数学处理为现代天体物理学提供了众多关于大质量恒星及黑洞晚期演化阶段的理论模型。[6][7]许多重要的科学概念、机构与发明均以他命名，包括著名的钱德拉塞卡极限与钱德拉X射线天文台。[8]

钱德拉塞卡出生于英属印度晚期（英属拉杰时期），一生研究范围极广，涉及恒星结构、白矮星、恒星动力学、随机过程、辐射输运、氢负离子的量子理论、流体力学与磁流体稳定性、湍流、椭球平衡体的平衡与稳定性、广义相对论、黑洞的数学理论以及引力波碰撞理论等多个物理领域。[9]在剑桥大学期间，他提出了一个解释白矮星结构的理论模型，首次将电子简并物质中电子速度导致的相对论性质量变化纳入计算。通过此理论，钱德拉塞卡得出结论：白矮星的质量不能超过太阳质量的1.44倍——这一极限即著名的钱德拉塞卡极限。钱德拉塞卡还修正了扬·奥尔特等人最早提出的恒星动力学模型，将银河系中恒星围绕银心旋转时受到的引力场涨落影响纳入考量。他为这一复杂的动力学问题给出了数学解——涉及二十个偏微分方程，并由此提出了一个新的物理量：动力摩擦。这种效应一方面会减缓恒星的运动速度，另一方面又有助于维持星团的整体稳定。钱德拉塞卡进一步将分析扩展至星际介质，证明银河中的气体与尘埃云分布极为不均匀，揭示了星系结构与动力学之间更深层次的联系。
\subsection{早年生活与教育}
钱德拉塞卡于**1910年10月19日**出生在**英属印度的拉合尔（Lahore，今属巴基斯坦）**，出身于一个**泰米尔人家庭**。[11] 他的父母分别是**希塔·巴拉克里希南（Sita Balakrishnan，1891–1931）**与**钱德拉塞卡拉·苏布拉马尼亚·艾耶尔（Chandrasekhara Subrahmanya Ayyar，1885–1960）**。[12] 钱德拉塞卡出生时，其父在**英属印度西北铁路局（Northwestern Railways）**担任**副审计长（Deputy Auditor General）**，因此一家居住在拉合尔。钱德拉塞卡（常被称为 **Chandra**）有两位姐姐——**拉贾拉克希米（Rajalakshmi）**与**巴拉帕尔瓦蒂（Balaparvathi）**；三位弟弟——**维什瓦纳坦（Vishwanathan）**、**巴拉克里希南（Balakrishnan）**与**拉马纳坦（Ramanathan）**；以及四位妹妹——**萨拉达（Sarada）**、**维迪娅（Vidya）**、**萨维特丽（Savitri）**和**森达丽（Sundari）**。他的**叔父**正是印度著名物理学家、诺贝尔奖得主**钱德拉塞卡拉·文卡塔·拉曼（Chandrasekhara Venkata Raman）**。他的母亲是一位热爱学术与文学的女性，曾将**易卜生（Henrik Ibsen）**的戏剧《**玩偶之家（A Doll’s House）**》翻译成泰米尔语，被认为是**启发钱德拉塞卡早年智识兴趣**的重要人物。[13]1916年，钱德拉塞卡一家从拉合尔迁往**阿拉哈巴德（Allahabad）**，并于**1918年定居在马德拉斯（Madras，今金奈）**。

钱德拉塞卡在**12岁之前一直由家中教师教授课程**。[13] 在中学阶段，他的父亲教授他**数学与物理**，而母亲则教他**泰米尔语**。1922年至1925年间，他就读于**马德拉斯特里普利坎的印度中学（Hindu High School, Triplicane, Madras）**。随后，他进入**马德拉斯总统学院（Presidency College, Madras）**就读，该校隶属于**马德拉斯大学（University of Madras）**。1925年至1930年间，他在此完成本科课程。1929年，他受到**阿诺德·索末菲（Arnold Sommerfeld）**一次讲座的启发，撰写了自己的第一篇论文——《**康普顿散射与新统计（The Compton Scattering and the New Statistics）**》。[14]1930年6月，他以优异成绩获得**物理学荣誉学士学位（BSc Hons. in Physics）**。同年7月，钱德拉塞卡获得**印度政府奖学金（Government of India Scholarship）**，前往**剑桥大学（University of Cambridge）**攻读研究生，并被**三一学院（Trinity College）**录取——这一机会由他曾通信交流的物理学家**R. H. 福勒（R. H. Fowler）**促成。在赴英旅途中，钱德拉塞卡利用航行时间推导出**白矮星中简并电子气体的统计力学理论**，并在福勒此前研究的基础上，加入了**相对论修正项**（详见后文“学术遗产”部分）。
